% ch3_spark.tex
\chapter{Apache Spark}
\label{ch:spark}

La thèse de référence présente d'abord le modèle d'exécution d'Apache Flink
avant de décrire l'implémentation de StreamKM++. Nous suivons la même logique
avec Apache Spark. Ce chapitre présente uniquement les mécanismes du framework
qui sont mobilisés par notre solution : les RDD et DataFrames, le traitement
par micro-batches de Structured Streaming, la diffusion de données avec
\texttt{broadcast}, la réduction arborescente et le partitionnement.

L'objectif n'est pas de proposer une présentation générale de Spark, mais de
mettre en évidence les propriétés qui influencent directement le portage de
StreamKM++. Leur utilisation concrète dans le projet est décrite au
chapitre~\ref{ch:implementation}.


% ---------------------------------------------------------------------------
\section{RDD et DataFrames}
% ---------------------------------------------------------------------------

Apache Spark fournit plusieurs abstractions pour manipuler des données
distribuées. Deux d'entre elles sont utilisées dans notre implémentation :
les \textit{Resilient Distributed Datasets} (RDD) et les DataFrames.

Un RDD est une collection distribuée et partitionnée d'objets. Les opérations
appliquées à un RDD, telles que \texttt{map}, \texttt{filter} ou
\texttt{mapPartitions}, sont exécutées indépendamment sur ses différentes
partitions lorsque cela est possible. Cette abstraction convient au traitement
des objets \texttt{Point} utilisés par StreamKM++, car l'algorithme repose
principalement sur des opérations impératives locales : calcul de distances,
construction de coresets et mise à jour des structures merge-and-reduce.

Les DataFrames représentent quant à eux des données distribuées organisées
selon un schéma. Spark peut optimiser les opérations qui leur sont appliquées
à travers son moteur d'exécution SQL, notamment lorsque les traitements
reposent sur des expressions, des projections ou des agrégations structurées.

Dans notre architecture, les DataFrames sont principalement utilisés aux
frontières du traitement, notamment avec Structured Streaming, tandis que la
logique algorithmique de StreamKM++ est exécutée sur des RDD de
\texttt{Point}. Cette séparation permet de conserver les composants
algorithmiques indépendants de Spark tout en exploitant les mécanismes de
distribution du framework.


% ---------------------------------------------------------------------------
\section{Traitement par micro-batches avec Structured Streaming}
% ---------------------------------------------------------------------------

Apache Spark Structured Streaming permet de traiter des flux de données à
l'aide de la même API que les traitements distribués classiques. Dans notre
implémentation, il est utilisé en mode \textit{micro-batch}.

Le flux d'entrée est alors découpé en lots successifs. Chaque micro-batch est
considéré comme un ensemble de données borné et peut être traité au moyen des
opérations Spark habituelles. La méthode \texttt{foreachBatch} permet
notamment d'appliquer une fonction spécifique à chaque lot produit par le
moteur.

Ce modèle diffère de celui de Flink utilisé dans la thèse, où le traitement
est organisé autour d'opérateurs de streaming continus et d'un état maintenu
au niveau des subtasks. Dans l'implémentation étudiée par Bitsakis, certaines
frontières du traitement doivent notamment être détectées au moyen de
mécanismes propres à Flink, tels que les watermarks.

Avec Structured Streaming, la frontière entre deux micro-batches est
explicitement fournie par le moteur d'exécution. Cette caractéristique
simplifie certaines étapes du portage, puisqu'un lot peut être traité comme
un ensemble borné avant que le résultat intermédiaire soit intégré à l'état
global de l'application.

Cette approche introduit toutefois une différence importante avec un
traitement événement par événement : les observations ne sont pas
nécessairement traitées immédiatement lors de leur arrivée, mais lors de
l'exécution du micro-batch auquel elles appartiennent. Ce compromis est
discuté plus en détail au chapitre~\ref{ch:discussion}.


% ---------------------------------------------------------------------------
\section{Diffusion des centres avec \texttt{broadcast}}
% ---------------------------------------------------------------------------

Certaines opérations nécessitent qu'une même donnée soit accessible à
l'ensemble des tâches Spark. C'est notamment le cas lors de l'évaluation de
la qualité d'un clustering : les centres calculés doivent être utilisés par
toutes les partitions contenant les observations originales.

Spark fournit pour cela le mécanisme de \textit{broadcast variables}. Une
valeur diffusée est transmise aux exécuteurs puis réutilisée localement par
les tâches qui en ont besoin, ce qui évite de sérialiser et de transférer la
même information à chaque opération.

Dans notre projet, ce mécanisme est utilisé par
\texttt{CostEvaluator.cost}. Les centres du clustering sont diffusés, puis
chaque partition calcule localement sa contribution à la somme des erreurs
quadratiques. Les résultats partiels sont ensuite agrégés afin d'obtenir le
coût global.

Cette organisation correspond au rôle joué dans la thèse par la diffusion des
centres vers les opérateurs chargés de calculer le coût sur les données
distribuées. L'utilisation concrète de cette opération est détaillée dans la
section~\ref{sec:impl-evaluation}.


% ---------------------------------------------------------------------------
\section{Réduction arborescente}
\label{sec:spark-treereduce}
% ---------------------------------------------------------------------------

La construction distribuée de StreamKM++ produit plusieurs résumés
intermédiaires, généralement un par partition. Ces résultats doivent ensuite
être combinés pour former un coreset global.

Une première approche consisterait à transférer tous les résumés vers le
driver puis à les fusionner séquentiellement. Cette solution concentre le
travail de réduction sur un seul processus et devient moins adaptée lorsque
le nombre de partitions augmente.

Spark fournit des opérations de réduction arborescente telles que
\texttt{treeReduce} et \texttt{treeAggregate}. Au lieu d'effectuer toutes les
combinaisons directement sur le driver, les résultats intermédiaires sont
fusionnés progressivement à plusieurs niveaux. Une partie du travail peut
ainsi être réalisée en parallèle avant que le résultat final ne soit renvoyé
au driver.

Ce mécanisme est particulièrement adapté au merge-and-reduce de StreamKM++,
car la fusion de deux résumés produit elle-même un résumé qui peut être
fusionné à un niveau supérieur.

Dans notre implémentation, \texttt{treeReduce} est utilisé pour combiner les
structures \texttt{BucketManager} construites indépendamment sur les
différentes partitions. Ce choix permet de conserver une graine propre à
chaque partition lors de la construction des coresets locaux et d'éviter
l'introduction d'un accumulateur initial partagé entre les partitions.

L'intérêt de cette réduction arborescente est principalement architectural :
elle évite de centraliser immédiatement toutes les opérations de fusion sur le
driver. Son impact réel sur le temps d'exécution doit toutefois être mesuré
expérimentalement et sera étudié dans la section consacrée à la scalabilité.


% ---------------------------------------------------------------------------
\section{Partitions et ressources de calcul}
\label{sec:spark-partitions-coeurs}
% ---------------------------------------------------------------------------

Le partitionnement constitue un élément central du modèle d'exécution Spark.
Un RDD est divisé en plusieurs partitions, chacune pouvant être traitée
indépendamment par une tâche.

Dans notre cas, chaque partition construit son propre
\texttt{BucketManager} et produit donc un résumé local des observations qui
lui sont attribuées. Le nombre de partitions influence ainsi directement la
structure de la construction distribuée : modifier ce nombre modifie le nombre
de coresets intermédiaires ainsi que la manière dont ils sont ensuite
fusionnés.

Il est important de distinguer ce paramètre du nombre de ressources de calcul
disponibles. Le nombre de partitions définit le nombre potentiel de tâches et
la décomposition logique des données, tandis que le nombre de c\oe{}urs
disponibles détermine combien de ces tâches peuvent effectivement être
exécutées simultanément.

Les deux paramètres ne doivent donc pas être confondus. Une augmentation du
nombre de partitions n'implique pas automatiquement une amélioration des
performances : un partitionnement plus fin peut augmenter les possibilités de
parallélisme, mais également introduire davantage de tâches, de sérialisation
et d'opérations de fusion.

De la même manière, le partitionnement peut influencer le coreset final,
puisque chaque partition réalise sa propre construction locale avec une graine
déterminée. L'effet sur la qualité du clustering n'est donc pas supposé a
priori ; il doit être observé expérimentalement.

Cette séparation entre partitionnement des données et ressources de calcul
constitue un point intéressant du portage vers Spark. Elle permet d'étudier
séparément l'effet du découpage des données et celui du parallélisme matériel,
ce qui sera exploité dans les expériences de scalabilité du
chapitre~\ref{ch:evaluation}.